\section{Introduction}\label{sec:intro}

Much of fundamental physics is organized around the hope of a final theory: a single
theory of everything from which every regularity follows, and which is not itself
the approximation of something deeper \citep{weinberg1992}. Some go further and hold
that the world simply \emph{is} such a structure \citep{tegmark2008}, or that a unique
computational object underlies all possible physics \citep{wolfram-ruliad2021}. We do
not ask whether a final theory exists. We ask whether one could ever be \emph{known} to
be final (\cref{def:known}), and we argue that it could not.

The argument starts from two premises. The first is nearly definitional: every
observation, measurement or other use of a system leaves a finite record written at
finite precision (\cref{prem:record}). The second is an assumption we assert and defend
rather than prove: every system anyone can use, including every method of choosing
among theories, is computable or a computable stochastic process (\cref{prem:effective}).
It is the modest physical Church--Turing thesis \citep{piccinini2011}, extended from
devices to methods.

From these premises the main claim follows in three parts.
\begin{enumerate}[label=(C\arabic*)]
  \item \emph{No certification.} Any theory that fits every observation made so far is one
        of infinitely many that fit them equally well and disagree about what has not yet
        been observed. So no method of any kind, computable or not, can certify a theory
        as final: if a method certifies a model, deterministic or stochastic, at some
        record, a computable world that produces that record makes the certification
        wrong (\cref{thm:k1-certify}).
  \item \emph{No reliable discovery.} No usable method is guaranteed, even with any
        fixed probability, to settle on a final theory. Whatever has been observed so far,
        there is a possible world that fits all of it in which a deterministic usable
        method never settles on the right theory; for a randomized usable method and
        every $p>0$, there is such a world in which it settles with probability less than
        $p$ (\cref{thm:ud-gold,ud:rem-randomized}). A user can therefore never rule out living in such a
        world. This part rests on the second premise; for deterministic methods it needs
        only the weaker premise that every world computable at a method's own level of
        power is possible (\cref{rem:p2star,thm:ud-gold-rel}).
  \item \emph{Permanent underdetermination, even in principle.} Even an ideal agent with
        unlimited data and computing power faces two limits. (a) Formulations that make
        identical predictions cannot be separated by any method. (b) Chance hypotheses
        that agree on which events have probability zero cannot be told apart
        \emph{with certainty} by any method (\cref{thm:t5-permanent}).
\end{enumerate}
A corollary follows (C$'$): simplicity narrows the field but cannot decide among the
survivors that matter. The data fix which models fit; which of them is shortest depends
on the reference machine. Rivals that fit equally well, are within $\K(x)+O(1)$ bits of
the shortest, and disagree at an unobserved setting $x$ always exist, and a change of
machine that moves no description length by more than about that gap makes such a rival the
unique shortest (\cref{thm:con-nearmin,thm:con-gap,con:cor:promote}). The universal
mixture predicts within $\K(\nu)$ bits of every model $\nu$ (\cref{thm:pred-relative}).
Every choice among those survivors---by
simplicity, beauty, naturalness or unification---is a constraint we impose
(\cref{lem:bias}). None of this makes the search for compact,
unified theories pointless. Compact theories compress what is known and commit to
predictions about what is not, which can be tested (\cref{sec:predict}). What is
misguided is the goal of a final theory known to be final.

\paragraph{Algorithmic agnosticism.}
We call this position \emph{algorithmic agnosticism}.\footnote{The name is unrelated
to ``agnostic learning'' in computational learning theory \citep{kearns1994}.} Huxley, who coined the word, stated agnosticism as a method:
do not pretend that conclusions are certain which are not demonstrated or demonstrable
\citep[pp.~245--246]{huxley1889}. In the sense we intend, commonly glossed as neither affirming
nor denying a claim that evidence cannot settle, we neither affirm nor deny that a final
theory exists; we hold that none could be known to be one. The position differs from van
Fraassen's constructive empiricism \citep{vanfraassen1980}, which is agnostic about
unobservable entities: we are agnostic about which law generates the observations
themselves. The qualifier \emph{algorithmic} names the framework of the argument, in
which theories are programs and usable methods are algorithms. For the contrast with
Müller's \emph{algorithmic idealism} \citep{muller2020,muller2026}, see
\cref{sec:objections}.

\paragraph{What we do not claim.}
We do not claim that the world is computable, discrete, continuous, random or
deterministic, nor that it lacks a true description.

\paragraph{Precursors.}
That universal laws can in principle be refuted but never verified goes back to Hume
and was made central by Popper \citep{popper1959}. That a finite record admits arbitrarily many laws,
so that a law is worth having only if it compresses its data, goes back to Leibniz and
Weyl \citep{leibniz1686,weyl1932} and was made precise by Chaitin \citep{chaitin2005}.
Putnam proved that no computable method extrapolates every computable sequence
\citep{putnam1963}, and Gold that none identifies every computable law in the limit
\citep{gold1967}; Kelly recast scientific method in these terms
\citep{kelly1996}. Assuming an objective, computable world, Hutter formalized theories
of everything as programs selected by description length, and called such a theory
``subjective'' because it must include the observer's localization within the world
\citep[Sec.~4]{hutter2010}. He also notes that ``for every model we can construct an
alternative model that behaves identically in the past'' and predicts differently
\citep[Sec.~4]{hutter2010}, which is our \cref{thm:ud-infinite}, and he formalizes and examines a
universal theory of everything that generates all computable universes
\citep[Sec.~6]{hutter2010}; choosing among such theories by description length is
machine-relative (\cref{thm:con-machine}). Schmidhuber treats a theory of everything as a
prior over computable universes and notes that assuming it formally describable
``imposes a strong inductive bias'' \citep[abstract]{schmidhuber2000}; \cref{lem:bias} makes the
point general. Sterkenburg argues that Solomonoff prediction cannot justify a
preference for simplicity: its predictors are mixtures over effective hypotheses, so
the operative assumption is effectiveness, and the resulting ``simplicity'' depends on
an arbitrary choice of universal machine \citep{sterkenburg2016}. A strong case for
simplicity is Vitányi's: ideal MDL is optimal when the data are typical for the true
hypothesis, and algorithmic statistics separates a record's meaningful information from
its noise \citep{vitanyi-li2000,gacs2001}; we answer it in \cref{sec:simplicity,sec:objections}.
Hawking argued from
G\"odel's theorem that a final theory may be out of reach \citep{hawking2002}, and
Faizal and colleagues argue from incompleteness that a theory of everything would require
non-algorithmic understanding \citep{faizal2025}. Kiefer concludes from G\"odel's theorems that,
unless space-time is fundamentally discrete, ``we can never decide whether a given
theory is the final one or not'' \citep[abstract]{kiefer2024}. Our route to a similar conclusion
runs through underdetermination by finite records and the limits of learning, not
through undecidability within the theory, so it does not depend on the theory's
arithmetic strength or on whether space-time is discrete. Hossenfelder argues that the
appeal to beauty and naturalness has misled the search for a final theory \citep{hossenfelder2018}.

\paragraph{Contributions.}
Little of the mathematics is new; the contribution is the assembly, aimed at a specific
target. Specifically:
(i) we show that the core of the claim, no certification (C1) and permanent
underdetermination (C3), holds for every method whatsoever;
(ii) we state Gold's theorem in a form in which, whatever record has been observed, some
world that fits it is one the method never gets right, which answers the reply that a
method need only succeed in \emph{our} world;
(iii) we make the machine-relativity of simplicity quantitative for theories fitted to
finite data (\cref{thm:con-machine}), rather than leaving it as a remark about arbitrary
constants;
(iv) we isolate the underdetermination that is permanent even for an ideal agent
(\cref{thm:t5-permanent}); and
(v) we reach a conclusion like Kiefer's (above) from finite records and the limits of
learning alone.

\paragraph{Outline.}
\Cref{sec:framework} fixes definitions and states and defends the premises.
\Cref{sec:certify} shows that no record certifies a final theory, and
\cref{sec:discover} that no usable method reliably finds one.
\Cref{sec:simplicity} shows that simplicity does not single one out, and
\cref{sec:ideal} what even an ideal agent cannot know.
\Cref{sec:predict} explains what the search is nonetheless good for.
\Cref{sec:objections} answers objections, and \cref{sec:conclusion} concludes.
